\documentclass[a4paper,11pt]{article}

\usepackage[margin=2cm]{geometry}
\usepackage{fontspec}
\usepackage{polyglossia}
\setdefaultlanguage{russian}
\usepackage{titlesec}
\usepackage{hyperref}
\usepackage{xcolor}
\usepackage{parskip}
\usepackage{enumitem}

\setmainfont{FreeSerif}
\setsansfont{FreeSans}
\setmonofont{FreeMono}

\hypersetup{
  colorlinks=true,
  linkcolor=black,
  urlcolor=blue
}

\titleformat{\section}{\Large\bfseries\sffamily}{}{0em}{}[\vspace{-0.4em}\rule{\textwidth}{0.5pt}]
\titlespacing*{\section}{0pt}{1em}{0.5em}
\setlist[itemize]{left=0pt, nosep, label={--}}
\emergencystretch 1em

\begin{document}

\begin{center}
  {\Huge\bfseries\sffamily Фриев Давид Сергеевич}
  \vspace{0.2em}

  19 лет, Москва $\cdot$
  \href{https://github.com/koftamainee}{github.com/koftamainee} $\cdot$
  \href{https://t.me/koftamainee}{@koftamainee} $\cdot$
  \href{mailto:dev@koftamainee.ru}{dev@koftamainee.ru} $\cdot$
  +7~953~901~74~79
\end{center}

\section{Обо мне}

Разрабатываю на Go и C++, 3 курс ПМИ РГУ им. А. Н. Косыгина.
Пишу микросервисы на Go (REST API, NATS, PostgreSQL, Redis), работаю с Docker и Linux.
Ищу позицию Junior Go Developer / C++ Developer или стажировку в backend-разработке.

\section{Проекты}

\subsection*{\href{https://github.com/koftamainee/nimbus}{Nimbus}}
\hfill \textit{Go, Rust, gRPC}

\begin{itemize}
  \item Распределённый оркестратор контейнеров: управляющий кластер (Quorum) и рантайм контейнеров (Forge)
  \item Quorum --- control plane на Go с собственной реализацией Raft-консенсуса, WAL-логирование, gRPC API, KV-хранилище с транзакциями
  \item Forge --- контейнерный рантайм на Rust: изоляция через Linux namespaces/cgroups v2, распаковка образов
  \item Планировщик размещения, цикл согласования состояния для самовосстановления при отказах узлов
\end{itemize}

\subsection*{\href{https://github.com/koftamainee/search-engine}{Search Engine}}
\hfill \textit{Go, Docker}

\begin{itemize}
  \item Микросервисная поисковая система: веб-краулер, сервис индексации (BM25), REST API на Go
  \item Очереди задач через Redis, данные в PostgreSQL и Meilisearch
  \item Развёрнут в Yandex Cloud через Docker Compose по адресу \href{https://search.koftamainee.ru}{search.koftamainee.ru}
  \item Архитектура с учётом горизонтального масштабирования микросервисов
\end{itemize}

\subsection*{\href{https://github.com/koftamainee/mantle}{Mantle}}
\hfill \textit{C++, Vulkan}

\begin{itemize}
  \item Модульный игровой движок с Vulkan-рендерером и Frame Graph для автоматической синхронизации GPU
  \item Кастомная система управления памятью (arena, tlsf аллокаторы)
  \item Интеграция ECS-архитектуры (flecs) и физического движка (Jolt Physics)
  \item CI/CD с автоматической сборкой под Linux и Windows
\end{itemize}

\subsection*{\href{https://github.com/koftamainee/mint}{Mint} / \href{https://github.com/koftamainee/mintload}{Mintload}}
\hfill \textit{Rust, C}

\begin{itemize}
  \item Ассет-пайплайн для игрового движка: конвертация glTF/GLB в бинарные форматы для прямой загрузки на GPU
  \item Система сцен с текстовым (.mscn) и бинарным (.mscb) форматами, инкрементальная сборка
  \item Mintload --- библиотека на C для быстрой загрузки файлов через mmap (с C++ обёрткой)
\end{itemize}

\subsection*{Другие}

\href{https://github.com/koftamainee/chalkboard}{Chalkboard} --- C++20 библиотека для рендеринга LaTeX-формул с hot-reload.

\href{https://github.com/koftamainee/danmaku}{Danmaku} --- bullet hell игровой движок на чистом C с Lua-скриптингом (SDL3).

\section{Навыки}

\begin{tabular}{@{}ll}
\textbf{Языки:} & C, C++, Go, Rust, Lua, Python \\
\textbf{Инструменты:} & Git, CMake, Docker, Docker Compose, Linux, CI/CD \\
\textbf{Базы данных:} & PostgreSQL, Redis, NATS \\
\textbf{Прочее:} & Vulkan, Nginx, REST API \\
\textbf{Английский:} & B2 (Upper-Intermediate) \\
\end{tabular}

\section{Образование}

  \textbf{РГУ им. А.Н. Косыгина}, Москва \\
  Прикладная математика и информатика, 3 курс (2024--2028)

\end{document}
